% !TeX root = ../papers/nc-p3.tex
\section{一つの多重点を持つ被約四次曲線の境界}
\label{sec:additive-pencil-boundary}
二つの接する枝、平面通常三重点、四枝の楕円多重点を同じ計算で扱う。
$C$の正規化を$m=2,3,4$本の射影直線とし、各枝の原点の局所座標を$z_i$とする。
それらの点を一つに合わせた局所環は、値が一致し、一次係数の和が零となる
正規化上の関数全体とする。枝のパラメータを変えれば二枝の場合は接する二枝の一次係数の一致となる。
直線束の貼り合わせでは、適切な自明化で、共通値$c$と一次係数$v_i$への条件は
$\sum v_i=\lambda c$となる。各枝での定数倍を除けば、この一つの加法パラメータが残る。
二つの束の積ではパラメータは加算される。

\begin{theorem}[正則性と被約境界の被覆]\label{thm:additive-degree-balance}
各成分上で正の次数を持ち総次数4の可逆層$L,M$と、$LM$の基点自由鉛筆を用いる。
定理\ref{thm:elliptic-pencil-all-as}の反転二窓が標準Koszul AS正則となる必要十分条件は
多重次数が一致することである。一致して$L\not\simeq M$なら反対階数は8で、
両側の二次核はそれぞれの完全線型系が定める$C$を切り出す。
添字保持同型の必要十分条件は、曲線同型が順序付きの二束と鉛筆を対応させることである。
このとき固定した多重次数での有理商の次元は$13-m$である。
\end{theorem}
\begin{proof}
導手は各枝の二重原点である。双対化層は自明である。実際各枝の$dz_i/z_i^2$を組にすると、
正則関数を掛けた留数和が零となる。逆に任意の双対化微分は導手条件で高々二重極となり、
各枝の二重極係数は等しく、単純極係数の和は零である。この組は上の微分の正則関数倍で尽くされる。
非負多重次数で総次数が正なら逆束の切断は零なので$H^1=0$である。
総次数が2以上なら元の束は大域生成される。共通値を非零に指定し、
正次数の枝の一次係数を調整して特異点で生成できる。
次数2以上の枝の他の点では導手で消える多項式で生成できる。次数1の枝ではもう一つの正次数の枝があれば
一次係数を互いに調整できる。正次数の枝が一つだけならその次数は2以上である。
次数零の枝は共通値で生成される。従って多重次数が一致する場合、
命題\ref{prop:elliptic-pencil-polygons}と同じ評価核のフィルトレーション、中心正則列、完全対合の
証明により全次数でAS正則性が従う。異なる場合は任意の真部分集合$I$上だけに支えを持つ
$L$の切断の次元が$d(I)-1$である。各枝で定数項を零にし、一次係数の和を零にするためである。
系\ref{cor:elliptic-polygon-degree-balance}と同じ相補的部分集合で指数成長の商を得る。

$|L|$は各枝を有理正規曲線に写す。線型包の座標は共通値、和に一式を課した一次係数、
枝ごとに独立な二次以上の係数である。したがって枝は原点以外で交わらず、局所環も復元される。
詳しくは、一次係数の超平面は埋込みの座標で実現する。二次以上を一枝だけで実現するには、
次数2以上の枝には二次の座標がある。$m\ge3$では指定枝と異なる一枝のみに
一次係数を持つ二関数を、異なる相手枝を選んで掛ければ指定枝だけの二次初項を得る。
$m=2$で一方が直線でも、他方の二次座標を共通接線座標の平方から引けばよい。
その後は一次座標を掛けて全ての高次初項を得る。完備化で導手全体を復元できるため閉埋込みとなる。
なお二枝の線型包は共通接線を共有するが、共通接線は非線型な有理正規曲線と二重原点でのみ交わる。
三枝以上の二つの線型包の交わりは共通点だけである。
この座標表示で平面に含まれる次数3の部分曲線はない。例えば三枝では二次曲線の内部座標と
独立な二つの接方向があり、四枝なら任意の三接方向は独立である。
従って前節と同じ次数の議論で二二次式完全交叉となる。
多重次数零の非自明な束では定数の組が一次貼り合わせを満たさず切断を持たない。
これで階数8と二次核が従う。完全交叉を平滑化して二束と鉛筆を延ばす議論も前節の通りである。

グラフの復元を詳しく示す。$r_i=2d_i$とし、鉛筆の一般の切断の零点を取ると、
前節と同様にテンソルは各枝の$r_i$個の異なる非特異点での非零重み付き評価の和である。
一枝だけに支えを持つ$L^2$の切断は、その枝で二重原点を引いた全切断である。
そのテンソル縮約の階数は$r_i-1$、像は対応する枝の評価の包に入る。
$d_i\ge2$ではこの階数は異なる枝の線型包の共通部分の次元より大きい。
$m\ge3,d_i=1$では像の生成元は二次Veronese上の異なる二点評価の非零重み付き和で、
他枝との共通点である原点評価の倍数にはならない。二枝で小次数の枝があれば、
大次数の枝の対応からその相手も定まる。

各正規化の原点を無限遠に置き、一方の候補グラフで両側のアフィン座標を$t$に揃える。
他の候補の枝写像を$t\mapsto at+b_i$と書く。傾き$a$は全枝で共通である。
これは二つの同定の比が、関数の一次係数の和零の条件を保つからである。
テンソルの$i$成分の表示は、$\deg u,\deg v\le r_i$に対して一方は
$\ell_i(u(t)v(t))$、他方は$\ell'_i(u(t)v(at+b_i))$となる（自明化の定数は汎関数に含める）。
$d_i\ge2$なら、片側の成分支持に制限した階数$r_i-1$のテンソルは、両候補の
$(r_i-2,r_i)$グラフの線型包に入る。異なる二グラフは長さ2で交わり、
合併のイデアルでひねると$\mathcal O(r_i-4,r_i-2)$で$H^1=0$になる。
よって二つの包の交わりの行列階数は高々2であり、その枝で同定は一致する。

$r_i=2$なら、源または先で次数0の多項式を使うとその切断は当該枝だけに支えを持つ。
相手側はこの枝への制限が全射なので、$m_j=\ell_i(t^j)$と置けば
$\ell_i'(t^j)=m_j$ ($0\le j\le2$)、
\[
 (a-1)m_1+b_i m_0=0,\qquad
 (a^2-1)m_2+2ab_i m_1+b_i^2 m_0=0.
\]
相異なる二点の重みは非零だから$m_0m_2-m_1^2\ne0$である。
大次数の枝があれば$a=1$となり、上の二式から$b_i=0$となる。

残るのは四直線の場合である。テンソルの二つの表示の差を枝ごとに取ると、
源か先の次数0の座標との組では差が零であり、各差は両側の原点の一次ジェットの包に入る。
全枝の和は零である。原点評価の直線で両側を割ると、残るのは接方向のテンソル
$h_i\otimes k_i$の係数だけである。三次元の四接方向には全係数非零の一つの関係だけがあり、
対応する四つの$h_i\otimes k_i$は独立である。実際、最初の三方向を両側の基底とすると、
第四項だけが基底の非対角積を持つ。従って各差で一次対一次の係数も零であり、
$m_2=a m_2+b_i m_1$を得る。最初の一次式と判別式の非消滅から$a=1,b_i=0$である。
全枝で一致すれば被約な曲線上の射は同じであり、鉛筆は隣接核から復元できる。

各多重次数での束の選択は一次元の加法群、鉛筆は12次元である。
曲線自己同型の連結部分は無限遠を保つ枝ごとのアフィン変換で、傾きが共通となるため次元$m+1$である。
一般の鉛筆の安定化群は有限である。各枝での二切断は一般に選べるので、次数4以上の枝があれば
その分岐値が三つ以上あり、基底の自己同型が有限に限られる。四直線のときも
異なる枝の二次射の分岐値の和集合は一般に三点以上となる。実際に二枝への切断の制限は
共通値を合わせれば一次係数の差を残る枝で調整でき、二つの二次射は無限遠での値だけを
共有して一般に選べる。基底変換を固定した後は各有限射の持ち上げは有限個である。
従って差は$14-(m+1)=13-m$となる。
\end{proof}

\begin{corollary}[被約二二次式完全交叉の曲線型]\label{cor:reduced-quartic-conductor}
$\mathbb P^3$内の被約二二次式完全交叉を曲線入力とし、各成分上で二束が正次数である鉛筆構成では、
上記と前節の曲線型で尽くされる。
\end{corollary}
\begin{proof}
$C$は連結で算術種数1、双対化層は自明である。正規化の成分の種数、成分と特異点の接続グラフの
第一Betti数、各特異点の$\delta_p-b_p+1$（$b_p$は枝数）の和は1である。
最後の数は、共通値だけを課す正規化の部分環からの余長であり非負である。
種数1の成分は導手因子の次数が$2-2g=0$なので他枝と交われず、滑らかな場合に尽きる。
ここで導手因子の次数は有限正規化の双対性
$\nu_*\omega_{\widetilde C}=\mathcal Hom_C(\nu_*\mathcal O_{\widetilde C},\omega_C)$から従う。

局所寄与が全て零なら、特異点の完備環は共通値だけを課す環である。
枝数$b$でその導手の正規化側の余長は$b$、$\delta=b-1$である。
Gorenstein曲線で導手のこの余長は$2\delta$である
（正規化完全列の双対を取り、長さ$\delta$の商の局所双対も同じ長さとなる）。
従って特異点は節点だけで、各有理成分上の導手の次数2から環状配置になる。自己節点の一成分も含む。
残りの場合は一つの特異点の局所寄与が1で、接続グラフは木となる。他の節点は分離節点になるが、
その一方の部分曲線は算術種数0であり、そこへの双対化層の制限次数は$-2+1=-1$で矛盾する。
従って他の特異点はなく、各有理成分は一枝だけでこの点と交わる。
各枝の導手次数は2なので、環は共通値に加えて一次係数の超平面を指定する。
この超平面の関係の全係数が非零でなければ、一枝の導手が一次となってしまう。
各枝の座標を変えて上の和零の表示を得る。一枝なら尖点であり、枝数は次数4以下である。
\end{proof}
